\section{Introduction}

Large language models (LLMs) have demonstrated strong reasoning ability on
tasks such as mathematical proof generation, complex question answering, and
code synthesis.
Yet their training and alignment are still dominated by outcome-level
supervision: rewards are typically assigned only to the final answer, while
intermediate steps in long reasoning chains receive little direct feedback.
This sparse-reward setting creates a severe temporal credit-assignment
problem.
A model can observe whether a trajectory eventually succeeds, but it cannot
directly tell which steps were contributed to  that success
\citep{uesato2022solving, lightman2023lets}.
As a result, process reward models (PRMs), which provide fine-grained
step-level feedback, have become an important component in test-time scaling,
reranking, and reinforcement learning for LLM reasoning.

The main challenge in developing current process reward models is not the
model architecture, but obtaining high-quality supervision for individual
reasoning steps.
Human annotation is reliable but expensive and difficult to scale.
Monte Carlo-based labeling is more automatic and therefore less expensive,
but it requires large rollout budgets and often remains noisy, since the
same step can receive different labels across sampled continuations and
value estimates.
LLM-as-a-judge further reduces annotation cost, yet its labels can inherit
biases from pretraining data and be sensitive to prompt wording or scoring
criteria, which makes the resulting supervision unstable.
More broadly, these approaches still rely on external annotators or
evaluators to determine which intermediate steps are valid.
This dependence limits robustness, scalability, and reusability.



Meanwhile, outcome-level reinforcement learning methods such as GRPO
optimize policies from terminal rewards or group-relative advantages,
which cannot distinguish pivotal reasoning steps from incidental
ones or locally plausible but globally unhelpful steps becausethis signal is still too coarse for step credit
assignment.
This challenge leads us to ask whether reliable step-level credit can be derived
directly from trajectory data itself, without relying on external judges.



We propose to formulate a multi-step reasoning trajectory as a cooperative
game, where reasoning steps are the players and the final trajectory reward is
the shared payoff.
The key insight of this view is to turn sparse outcome supervision into an
allocation problem: step credit assignment becomes the problem of distributing
the shared payoff back to individual reasoning steps.
Without a principled allocation rule, such credit is often assigned through
coarse heuristics such as reward broadcasting or position-based proxies.
The remaining question is how to obtain a principled allocation rule within
this cooperative game.
Shapley values provide a natural starting point because they allocate shared
payoff by marginal contribution, offering a way to turn sparse final rewards
into step-level credit.


However, reasoning trajectories differ fundamentally from the classical
Shapley setting.
Reasoning steps are temporally ordered, causally dependent, and not freely
exchangeable.
Under strict causal constraints,
naively applying classical Shapley can
degenerate into prefix-wise value differences along the original sequence,
which mainly redistributes the same sparse outcome signal and remains sensitive
to value-estimation noise.
As a result, it fails to capture non-local interactions across steps in a
usefully dense form.
Therefore, the key design problem is not merely to adopt Shapley values, but
to adapt their fairness intuition to ordered, causally constrained reasoning
processes.

Motivated by this goal, we propose \textbf{SPRM}
(\textbf{S}hapley-value based \textbf{P}rocess \textbf{R}eward
\textbf{M}odel), a framework that automatically constructs trainable
step-level supervision from outcome-only signals.
Our central idea is to keep the marginal-contribution intuition of Shapley
values while adapting it to ordered reasoning trajectories.
SPRM combines trajectory-specific local credit, which captures how a step
changes the value of its own reasoning chain, with a cross-trajectory global
prior, which captures whether a step pattern is broadly associated with
successful reasoning rather than one noisy rollout.
This yields step supervision that is both causally grounded within a
trajectory and more stable across trajectories.

This design directly addresses the limitations of existing PRM supervision
pipelines.
Unlike annotation-based\cite{uesato2022solving,lightman2023lets} or judge-based approaches\cite{zeng2025versaprm, duan2025efficient,yang2025beyond,wang2025athena,cao2025dreamprm}, SPRM derives supervision
directly from sparse outcome signals and therefore removes the need for human
step labels or external evaluators.
Unlike purely local counterfactual scoring, it combines continuous local credit
with a discrete counterfactual check and global semantic priors, making the
resulting labels more
stable and less sensitive to noise in any single trajectory.
Unlike outcome-level RL alone, SPRM produces reusable step-level reward
signals that can support PRM training, reranking, and test-time reasoning.
Overall, SPRM aims to provide a more coherent trade-off among automation,
causal grounding, and scalability for process supervision in LLM reasoning.

Our main contributions are summarized as follows:

\begin{itemize}
    \item We propose an annotation-free automatic labeling method that
    converts outcome-only rewards in LLM reasoning into trainable step-level
    supervision, without relying on human annotation or external judges.
    \item We formulate a multi-step reasoning trajectory as a shared-payoff,
    ordered, and causally constrained credit-assignment problem, and
    introduce a Shapley-inspired attribution principle tailored to this
    setting.
    \item We present an end-to-end SPRM training framework that jointly
    models trajectory-internal local credit and cross-trajectory global
    priors, and validate its effectiveness through comprehensive analyses
    and experiments on ProcessBench, PRMBench and Best-of-\(N\)
    benchmarks.
\end{itemize}